\documentclass[10pt,a4paper]{amsart}
\usepackage[margin=19mm]{geometry}
\usepackage{amsmath,amssymb,mathtools}
\usepackage[scr=rsfs,cal=euler]{mathalfa}
\usepackage{xfrac}
\usepackage[unicode,hidelinks,bookmarks=false]{hyperref}
\newcommand{\ie}{i.e.\ }
\newcommand{\cf}{cf.\ }
\newcommand{\resp}{respectively\ }
\newcommand{\Subsection}{\subsection}
\providecommand{\nobreakdash}{\nobreak\hbox{-}}
\setlength{\emergencystretch}{2em}
\multlinegap0pt
\usepackage{bm}
\usepackage{bbm}
\usepackage{dsfont}
\usepackage{xspace}
\usepackage{cancel}
\usepackage{mathtools}
\usepackage{amsthm}
\makeatletter
\@ifundefined{theorem}{\newtheorem{theorem}{Theorem}}{}
\@ifundefined{proposition}{\newtheorem{proposition}[theorem]{Proposition}}{}
\makeatother
\newcommand{\E}{\mathbb{E}}
\newcommand{\W}{\mathbf{W}}
\newcommand{\h}{\mathbf{h}}
\title[Optimal Depth of Neural Networks]{Optimal Depth of Neural Networks}
\author{Qian Qi}
\date{Source: arXiv:2506.16862v1 (CC BY 4.0)}
\begin{document}
\maketitle

\noindent\emph{Source and licence.} Optimal Depth of Neural Networks. Version arXiv:2506.16862v1 (CC BY 4.0), distributed under the Creative Commons Attribution 4.0 International licence, as recorded in the arXiv licence metadata for this exact version and shown on the abstract page https://arxiv.org/abs/2506.16862v1. Full rights notice: licenses/arxiv-2506.16862-CC-BY-4.0.txt.\par\smallskip

\noindent\textit{Editorial scope.} Counted unit: the Proposition whose source label is \texttt{None} in the article (source0.tex lines 626--631, source label \texttt{None}), delivered together with the article's own complete original proof of it (source0.tex lines 633--660, one \texttt{proof} environment of 28 lines). No mathematics is written, repaired or re-derived: statement and proof are the article's own text line for line. Required context retained uncounted: none. Every definition, display and label of those lines is retained unchanged. Omitted from this package and not redistributed: 1--625, 632--632, 661--713. Nothing inside a retained excerpt is omitted or altered: every retained body is delivered exactly as the article's lines are written, so no omission receipt of any kind accompanies this package. Citations: the retained excerpts carry no citation command at all, so no bibliography entry of the article is transcribed and no reference is redistributed. Reference adaptation: none was needed and none was made; every reference inside the delivered unit resolves inside it, so no unresolved reference or citation is produced. Printed slips kept literal: no printed word, symbol, hypothesis or number of the counted statement or of its proof was altered. Layout and class: the article's own class is \texttt{imsart} with packages including inputenc, fontenc, amsthm, amsmath, amssymb, amsfonts, graphicx, booktabs, lmodern, tikz, pgfplots, microtype, hyperref, natbib; the wrapper is the lane's pinned \texttt{amsart} editorial wrapper at 10pt on A4, which restates the theorem-like environments the retained text needs and adds the article's own macro definitions that the retained text needs (\texttt{\textbackslash E}, \texttt{\textbackslash W}, \texttt{\textbackslash h}), with extra wrapper packages (bm, bbm, dsfont, xspace, cancel, mathtools). The delivered numbering is the unit's own. Assets: no retained span contains a figure environment, a graphics-inclusion command, a PDF-page inclusion, a table or a TikZ picture; no image file is copied into this package, the package declares no asset dependency and no image file is redistributed with it. Licence: version arXiv:2506.16862v1 of this article is distributed under the Creative Commons Attribution 4.0 International licence, as recorded in the arXiv licence metadata for this exact version and shown on the abstract page https://arxiv.org/abs/2506.16862v1. A full rights notice is delivered at licenses/arxiv-2506.16862-CC-BY-4.0.txt. Characters: every retained excerpt is pure ASCII, so the delivered engine, XeLaTeX with its default Latin Modern text font, reports no missing character.
\begin{proposition}
Let the utility function $g$ be Lipschitz continuous with constant $K_g$ (per Assumption A2). Suppose there exists a layer index $L_0$ such that for all layers $l \ge L_0$, the expected norm of the residual function's output is bounded by some $\delta > 0$, i.e., $\E[\|f_l(\h_l)\|] \le \delta$. If the per-layer cost $c$ satisfies $c > K_g \cdot \delta$, then for all $l \ge L_0$, the expected one-step gain in utility is strictly less than the per-layer cost:
\begin{equation}
\E[g(\h_{l+1}) - g(\h_l)] < c.
\end{equation}
\end{proposition}

\begin{proof}
We want to bound the expected change in utility from layer $l$ to $l+1$. Let $l \ge L_0$.
\begin{enumerate}
    \item \textbf{Start with the quantity of interest.} The expected one-step gain in utility is $\E[g(\h_{l+1}) - g(\h_l)]$.

    \item \textbf{Apply the ResNet update rule.} Substitute $\h_{l+1} = \h_l + f_l(\h_l; \W_l)$:
    \[ \E[g(\h_{l+1}) - g(\h_l)] = \E[g(\h_l + f_l(\h_l)) - g(\h_l)]. \]

    \item \textbf{Use the property $x \le |x|$.}
    \[ \E[g(\h_l + f_l(\h_l)) - g(\h_l)] \le \E\left[ \left| g(\h_l + f_l(\h_l)) - g(\h_l) \right| \right]. \]

    \item \textbf{Apply Lipschitz continuity.} By Assumption A2, the utility function $g$ is Lipschitz continuous with constant $K_g$. This means that for any two points $\mathbf{a}, \mathbf{b}$ in its domain, $|g(\mathbf{a}) - g(\mathbf{b})| \le K_g \|\mathbf{a} - \mathbf{b}\|$. Applying this inside the expectation with $\mathbf{a} = \h_l + f_l(\h_l)$ and $\mathbf{b} = \h_l$:
    \begin{align*}
    \E\left[ \left| g(\h_l + f_l(\h_l)) - g(\h_l) \right| \right] &\le \E[K_g \cdot \|(\h_l + f_l(\h_l)) - \h_l\|] \\
    &= \E[K_g \cdot \|f_l(\h_l)\|].
    \end{align*}

    \item \textbf{Use linearity of expectation.} The constant $K_g$ can be pulled out of the expectation:
    \[ \E[K_g \cdot \|f_l(\h_l)\|] = K_g \cdot \E[\|f_l(\h_l)\|]. \]

    \item \textbf{Apply the proposition's hypothesis.} For $l \ge L_0$, we assumed that $\E[\|f_l(\h_l)\|] \le \delta$. Substituting this into our inequality chain gives:
    \[ \E[g(\h_{l+1}) - g(\h_l)] \le K_g \cdot \E[\|f_l(\h_l)\|] \le K_g \cdot \delta. \]

    \item \textbf{Apply the cost condition.} The final assumption is that the cost $c$ was chosen such that $c > K_g \cdot \delta$. Combining this with the previous step yields the final result:
    \[ \E[g(\h_{l+1}) - g(\h_l)] \le K_g \cdot \delta < c. \]
\end{enumerate}
Thus, for any layer $l \ge L_0$, the expected gain in utility is strictly smaller than the cost of computation for that layer.
\end{proof}

\end{document}
